\section{Conclusion}

\label{sec:conclusion}
% ============================================================================

The AI cloud industry is built on a broken trust model. Users send their most
sensitive data to providers who promise---but cannot prove---that they will not
read, store, or misuse it. This trust model fails under adversarial conditions
(breaches, subpoenas, rogue employees) and is incompatible with regulated
industries that require cryptographic guarantees.

Hanzo AI Cloud replaces trust with proof:

\begin{enumerate}[leftmargin=1.4em]
    \item \textbf{Confidential inference}: CKKS FHE ensures the provider
    evaluates models on encrypted inputs. The provider never sees the
    plaintext. This is a mathematical guarantee, not a policy.
    \item \textbf{Agent sovereignty}: Each agent's vault, memory, and identity
    are cryptographically isolated. The platform cannot access agent state
    without the agent's keys.
    \item \textbf{Verifiable serving}: Every inference result carries a
    receipt anchored to the Lux A-Chain, proving provenance. Regulators,
    auditors, and users can verify exactly which model produced which output.
    \item \textbf{Decentralized provision}: No single provider controls the
    infrastructure. Stake-weighted selection and on-chain SLA enforcement
    align provider incentives with user outcomes.
    \item \textbf{Edge-to-cloud continuum}: The same model format, the same
    attestation, the same CRDT sync protocol work from a phone to a data
    center. Inference happens where it makes sense.
\end{enumerate}

The era of ``trust us, we won't look at your data'' is ending. The era of
``we cannot look at your data'' is beginning. Hanzo AI Cloud is the
infrastructure for the latter.

\paragraph{Settlement on the Hanzo AI Chain (AIVM).}
The confidential-inference, agent-hosting, and model-serving primitives
described in this paper settle on the Hanzo AI Chain. The AI Chain is
the Lux unified A-Chain (LP-134) re-framed as the AI Virtual Machine,
secured by Quasar 3.0 three-lane post-quantum certificates (LP-105),
executed on QuasarGPU (LP-132), and scheduled per QuasarSTM lanes
(LP-010). It mints \textbf{AI}, the native coin of the Hanzo Network,
and hosts the on-chain HMM registry, agent registry, and three composable
compute markets (inference, training, GPU pool). Confidential inference
is delegated to F-Chain (LP-013, LP-067) through the \texttt{F\_CALL}
precompile; threshold custody for model weights is delegated to M-Chain.
See the companion paper \emph{Hanzo AI Chain (AIVM)}
(\texttt{\textasciitilde/work/hanzo/papers/hanzo-ai-chain/}); activation
is 25 December 2025 alongside Lux Quasar 3.0.

% ============================================================================
% REFERENCES
% ============================================================================
\begin{thebibliography}{99}

\bibitem{hanzocandle}
Z.~Kelling, D.~Wei, and M.~Chen.
\newblock Hanzo Candle: A {Rust} {ML} framework for high-performance inference.
\newblock \emph{Hanzo AI Research}, August 2024.

\bibitem{hanzofhe}
Z.~Kelling.
\newblock Hanzo {FHE} Inference: Fully homomorphic encryption for confidential
{AI} inference.
\newblock \emph{Hanzo AI Research}, April 2026.

\bibitem{hanzojin}
Z.~Kelling.
\newblock Jin: Unified multimodal {AI} architecture for vision-language, audio,
and cross-modal understanding.
\newblock \emph{Hanzo AI Research}, April 2023.

\bibitem{hanzoedge}
Z.~Kelling.
\newblock Hanzo Edge: On-device {AI} inference for mobile and {IoT} deployments.
\newblock \emph{Hanzo AI Research}, July 2024.

\bibitem{hanzobaseruntime}
Z.~Kelling.
\newblock Hanzo Base: The single-binary runtime for sovereign applications.
\newblock \emph{Hanzo AI Research}, April 2026.

\bibitem{luxwhitepaper}
{Lux Partners}.
\newblock Lux: A scalable multi-consensus blockchain with post-quantum
cryptography.
\newblock \emph{Lux Network Technical Report}, 2024.

\bibitem{hanzoplatform}
M.~Chen, D.~Wei, and Z.~Kelling.
\newblock Hanzo Platform: {PaaS} for {AI}-native application deployment.
\newblock \emph{Hanzo AI Research}, December 2021.

\bibitem{ckks2017}
J.~H.~Cheon, A.~Kim, M.~Kim, and Y.~Song.
\newblock Homomorphic encryption for arithmetic of approximate numbers.
\newblock In \emph{ASIACRYPT}, Springer, 2017.

\bibitem{openaiapi}
{OpenAI}.
\newblock {API} data usage policies.
\newblock \url{https://openai.com/policies/api-data-usage}, 2024.

\bibitem{gentry2009}
C.~Gentry.
\newblock Fully homomorphic encryption using ideal lattices.
\newblock In \emph{STOC}, ACM, 2009.

\bibitem{brakerski2014}
Z.~Brakerski, C.~Gentry, and V.~Vaikuntanathan.
\newblock ({Leveled}) Fully homomorphic encryption without bootstrapping.
\newblock \emph{ACM Transactions on Computation Theory}, 6(3):1--36, 2014.

\bibitem{gilad2016}
R.~Gilad-Bachrach, N.~Dowlin, K.~Laine, K.~Lauter, M.~Naehrig, and
J.~Wernsing.
\newblock {CryptoNets}: Applying neural networks to encrypted data with high
throughput and accuracy.
\newblock In \emph{ICML}, 2016.

\bibitem{luxagent}
Z.~Kelling.
\newblock Lux Agent Sovereignty: Decentralized identity and capability
management for autonomous agents.
\newblock \emph{Lux Network Technical Report}, 2025.

\end{thebibliography}

